% !TEX root = tarasov-qasp-zh.tex
\cleardoublepage
\thispagestyle{empty}
\begin{figure*}[!ht]
\centering
\includegraphics[width=0.73\textwidth]{figs/sec/tarasov-qasp-sec-01.pdf}
\end{figure*}
\begin{fullwidth}
{\textsf{\Large \hlred{不要忽视运动学！物体在空间与时间中如何运动的问题，无论从物理的角度还是从实用的角度看，都是相当有趣的。}}}
\end{fullwidth}
\cleardoublepage


\chapter[你能分析直线运动的运动学图像吗？]{你能分析直线运动\\的运动学图像吗？\label{ch-01}}

\begin{dialogue}

\tchr 你们已经见过匀变速直线运动（straight-line, uniformly variable motion）中，物体的速度和路程随时间变化的图像。据此，我想提出下面的问题。\hlred{给你如图~\ref{fig-01} 所示的速度图像，请据此画出路程随时间变化的图像。}
\pnote{术语：原书的"路程"（distance travelled）相当于人教版中的路程；人教版 $x$--$t$ 图像中的 $x$ 则是位移。本章各段时间均为正、运动方向不变，路程与位移大小相等；一旦往返，两者就不再相等。}

\std 可是这种图像我从来没画过。

\tchr 这不会有困难。不过，我们一起来推理。首先把整个时间分成三段：1、2、3（见图~\ref{fig-01}）。在第 1 段时间内物体怎样运动？这段路程的公式是什么？

\begin{figure}[!h]%
\centering
\includegraphics[width=0.6\textwidth]{figs/ch-01/fig-001.pdf}
\caption[根据这个图像，你能画出路程随时间变化的图像吗？]{根据这个图像，你能画出路程随时间变化的图像吗？\label{fig-01}}
\end{figure}

\std 在第 1 段，物体做初速度为零的匀加速运动。路程的公式是
\begin{equation}%
s(t) = \frac{at^2}{2}
\label{eq-001}
\end{equation}
其中 $a$ 是物体的加速度。

\tchr 利用速度图像，你能求出加速度吗？

\std 能。加速度是单位时间内速度的变化量。它等于线段 $\overline{AC}$ 与线段 $\overline{OC}$ 的长度之比。

\tchr 好。现在来看第 2 段和第 3 段。
\std 在第 2 段，物体以第 1 段末获得的速度 $v$ 做匀速运动。路程公式是 $s=vt$。

\tchr 慢着，你的回答不准确。你忘了匀速运动不是从初始时刻开始的，而是从时刻 $t_{1}$ 开始的。在这之前，物体已经走过了等于 $\dfrac{at_{1}^{2}}{2}$ 的路程。第 2 段路程随所历时间的变化由下式表示：
\begin{equation}%
s(t) = \frac{at_{1}^{2}}{2}+ v (t-t_{1})
\label{eq-002}
\end{equation}
记住这一点，请写出第 3 段的路程公式。

\std 第 3 段物体的运动是匀减速的。如果我没理解错，这段路程的公式应该是
$$
s(t) = \frac{at_{1}^{2}}{2}+ v (t_{2}-t_{1}) +  v(t-t_{2}) - \frac{a_{1}(t-t_{2})^2}{2}
$$
其中 $a_{1}$ 是第 3 段的加速度。它只有第 1 段加速度 $a$ 的一半，因为第 3 段的时间是第 1 段的两倍。

\tchr 你的式子还可以化简一些：
\begin{equation}%
s(t) = \frac{at_{1}^{2}}{2}+ v (t-t_{1}) - \frac{a_{1}(t-t_{2})^2}{2}
\label{eq-003}
\end{equation}
现在，把式~\eqref{eq-001}、\eqref{eq-002} 和~\eqref{eq-003} 的结果汇总起来吧。

\std 我明白了。路程图像在第 1 段是抛物线，第 2 段是直线，第 3 段是另一条抛物线（翻转过来，上凸、开口向下）。这是我画的图像（图~\ref{fig-02}）。

\begin{figure}%
\centering
\includegraphics[width=0.6\textwidth]{figs/ch-01/fig-002.pdf}
\caption[学生画错的路程 s 随时间 t 变化的图像。]{学生画错的路程 $s$ 随时间 $t$ 变化的图像。\label{fig-02}}
\end{figure}

\tchr 你的图有两处毛病：\hlblue{路程图像不应有折点（kink）。}它应当是光滑曲线，也就是说，抛物线应与直线相切。此外，\hlblue{上面那条（翻转的）抛物线的顶点应对应时刻 $t_{3}$。}正确的图像是这样的（图~\ref{fig-03}）。

\begin{figure}[!ht]%
\centering
\includegraphics[width=0.6\textwidth]{figs/ch-01/fig-003a.pdf}
\caption[更正后的路程 s 随时间 t 变化的图像。]{更正后的路程 $s$ 随时间 $t$ 变化的图像。\label{fig-03}}
\end{figure}

\std 请解释一下。 \\
\tchr 我们来看路程\emph{—}时间图像的一段（图~\ref{fig-04}）。\pnote{底本此图的图题误沿了图~\ref{fig-03} 的文字（"更正后的图像"），与本图内容（弦、切线与瞬时速度的确定）不符，显系录入之误；中文版按图面内容拟写图题，待考原书。}在从 $t$ 到 $t+\Delta t$ 的时间间隔内，物体的平均速度等于
\begin{equation*}%
\frac{s(t+\Delta t) - s(t)}{\Delta t}= \tan \alpha
\end{equation*}
其中 $\alpha$ 是弦（chord）$AB$ 与水平方向的夹角。要确定物体在时刻 $t$ 的速度，需要求 $\Delta t \rightarrow 0$ 时这些平均速度的极限。于是
\begin{equation}%
v(t)=\lim_{\Delta t \rightarrow 0} \frac{s(t+\Delta t)-s(t)}{\Delta t}
\label{eq-004}
\end{equation}
\pnote{式~\eqref{eq-004} 就是"瞬时速度等于路程对时间的导数"，与高中数学中一元函数导数的定义完全一致；下文"速度图像下方的面积等于路程"则是定积分的思想。中学物理不使用极限记号，但竞赛与大学普通物理以此为标准语言。}在极限过程中，弦变成经过点 $A$ 且与路程\emph{—}时间曲线相切的切线（tangent）（见图~\ref{fig-04} 中的虚线）。这条切线与水平方向夹角的正切，就是时刻 $t$ 的速度值。于是，由路程\emph{—}时间曲线上相应点处切线的倾角，就能求出任意时刻的速度。
\begin{figure}[!ht]%
\centering
\includegraphics[width=0.5\textwidth]{figs/ch-01/fig-004a.pdf}
\caption[路程—时间图像的一段：弦、切线与瞬时速度的确定。]{路程\emph{—}时间图像的一段：弦 $AB$、$A$ 点处的切线与瞬时速度的确定。\label{fig-04}}
\end{figure}
但是我们回到你画的图（见图~\ref{fig-02}）。由你的图像会得出：在时刻 $t_{1}$（以及 $t_{2}$），物体的速度有两个不同的值。从左侧趋近 $t_{1}$ 时，速度等于 $\tan \alpha_{1}$；从右侧趋近时，速度等于 $\tan \alpha_{2}$。按你的图像，物体在时刻 $t_{1}$（$t_{2}$ 同样）的速度必须发生突变（discontinuity），而实际上并没有（图~\ref{fig-01} 中的速度\emph{—}时间图像是连续的）。

\std 我懂了。\hlblue{速度图像的连续性，导致路程\emph{—}时间图像的光滑性。}
\pnote{物理上，位置不能突变（否则速度无穷大）；速度突变则要求冲击式的作用（碰撞、绳突然张紧），只有这时路程\emph{—}时间图像才会出现折点——这正是下面图~\ref{fig-05} 那道题的物理背景。至于 $t_{1}$、$t_{2}$ 处"光滑但曲率突变"，对应的是加速度突变：普通（非冲击）力作用下这是允许的。}

\tchr 顺便指出，两条抛物线的顶点应分别对应时刻 0 和 $t_{3}$，因为\hlblue{在这些时刻物体的速度为零，曲线在这些点的切线必须水平}。
\pnote{速度在 $t_{3}$ 减到零后，若物体就此停住（如刹车），路程\emph{—}时间图像此后是水平直线；把匀减速公式一直外推到停止之后得出"反向运动"，是中学习题里的经典错误（"刹车陷阱"）。}
\hlred{现在，利用图~\ref{fig-01} 的速度图像，求物体到时刻 $t_{2}$ 走过的路程。}

\std 先由速度图像定出第 1 段的加速度 $a$，再定出第 2 段的速度 $v$，然后利用公式~\eqref{eq-002}。物体在时间 $t_{2}$ 内走过的路程等于
\begin{equation*}%
s(t_{2}) = \frac{at_{1}^{2}}{2}+v(t_{2}-t_{1})
\end{equation*}
\pnote{底本此处误排为 $v(t_{2}-t)$，按公式~\eqref{eq-002} 应为 $v(t_{2}-t_{1})$，中文版径改。}

\tchr 完全正确。不过还有更简单的办法。物体在时间 $t_{2}$ 内走过的路程，在数值上等于速度\emph{—}时间图像下方、区间 $Ot_{2}$ 上图形 $OABD$ 的面积。
\pnote{"面积法"：人教版必修第一册第二章《匀变速直线运动的研究》正是用 $v$--$t$ 图像下的面积（分割、求和、取极限）推导匀变速位移公式的。注意若速度取负值，$t$ 轴下方的面积对应负位移；本章未涉及。}让我们再考虑一个问题，来巩固所学。\hlred{设路程\emph{—}时间图像有折点，如图~\ref{fig-05} 所示，其中曲线是以 $A$ 为顶点的抛物线。请画出速度\emph{—}时间图像。}

\begin{figure}[!h]%
\centering
\includegraphics[width=0.5\textwidth]{figs/ch-01/fig-005a.pdf}
\caption[带折点的路程 s 随时间 t 变化的图像。]{带折点的路程 $s$ 随时间 $t$ 变化的图像。\label{fig-05}}
\end{figure}

\std 既然路程图像有折点，那么在相应的时刻（$t_{1}$ 和 $t_{2}$），速度图像就应有突变。这是我的图像（图~\ref{fig-06}）。

\tchr 好。线段 $\overline{BC}$ 的长度是多少？
\begin{figure}[!h]%
\centering
\includegraphics[width=0.5\textwidth]{figs/ch-01/fig-006a.pdf}
\caption[与图 5 对应的速度 v 随时间 t 变化的图像。]{与图~\ref{fig-05} 对应的速度 $v$ 随时间 $t$ 变化的图像。\label{fig-06}}
\end{figure}
\std 它等于 $\tan \alpha_{1}$（见图~\ref{fig-05}）。不过角度 $\alpha_{1}$ 的数值我们并不知道。
\tchr 尽管如此，确定 $\overline{BC}$ 的长度不会有困难。注意：物体到时刻 $t_{3}$ 走过的路程，与它一直做匀速运动走过的路程相同（图~\ref{fig-05} 中从 $t_{2}$ 到 $t_{3}$ 的直线是从 0 到 $t_{1}$ 那段直线的延长线）。既然路程由速度图像下方的面积量度，可知图~\ref{fig-06} 中矩形 $ADEC$ 的面积等于三角形 $ABC$ 的面积。因此 $BC=2EC$，也就是说，从左侧趋近时刻 $t_{2}$ 时的速度，等于 0 到 $t_{1}$、$t_{2}$ 到 $t_{3}$ 两段匀速运动速度的两倍。

\end{dialogue}
